\Needspace{9\baselineskip}
\section{From continuous leaves to periodic invariant graphs}
\label{sec:geometry}

We first extract the phase-space information needed by the analytic
argument.  This step uses the continuity of the foliation, both twist
inequalities, and finite invariant area.  It does not assume that a leaf
is a graph or that it comes from a smooth caustic. Birkhoff's
invariant-curve graph theorem is classical for exact symplectic twist maps
\citep{Arnaud2009}; the direct argument here also records the rational
periodicity required by the subsequent stationary-action construction.

\subsection{Oriented lines and the two twists}

Translate the origin into $\Omega$.  Write
\[
 n_\theta=(\cos\theta,\sin\theta),\qquad
 e_\theta=(-\sin\theta,\cos\theta),\qquad
 h(\theta)=\max_{z\in\overline\Omega}z\cdot n_\theta.
\]
The normal-angle parametrization of the boundary is smooth: strict
convexity makes its normal map bijective, and positive curvature makes
its derivative nonzero.  Consequently
\begin{equation}\label{geo:support}
 \begin{gathered}
 h>0,\qquad r=h+h''>0,\\
 \gamma(\theta)=h(\theta)n_\theta+h'(\theta)e_\theta,
 \qquad \gamma'(\theta)=r(\theta)e_\theta.
 \end{gathered}
\end{equation}
Represent a billiard state by its outgoing oriented line, with velocity
$e_\theta$ and equation $z\cdot n_\theta=p$.  The starting point is the
backward intersection of that line with the boundary.  Both intersections
are transverse and uniquely determined, so this is a smooth change of
coordinates onto
\[
 \mathcal B=\{(\theta,p):\theta\in\T,
                  \ -h(\theta+\pi)<p<h(\theta)\}.
\]
We continue to denote the conjugated billiard map by $T$ and the given
annulus by $U$.

Choose the lift $\widetilde T$ for which the angular advance lies in
$(0,2\pi)$.  If $\gamma(\psi)$ is the forward collision point, reflection
gives
\begin{equation}\label{geo:reflection}
 \begin{gathered}
 \theta'=\theta+2a,\qquad \psi=\theta+a,\qquad 0<a<\pi,\\
 p=h(\psi)\cos a-h'(\psi)\sin a,\qquad
 p'=h(\psi)\cos a+h'(\psi)\sin a.
 \end{gathered}
\end{equation}
Indeed $e_\theta\cdot n_\psi=\sin a>0$, and reflection changes
$e_\theta$ to $e_{\theta+2a}$.  At fixed $\theta$, differentiation of
$p$ with respect to $a$ gives $-r(\theta+a)\sin a<0$.
At fixed $\theta'$, write $\psi=\theta'-a$; differentiation of $p'$
again gives $-r(\theta'-a)\sin a<0$.  Writing $\theta_{-1}$ for the
preceding lifted angle, we obtain
\begin{equation}\label{geo:twists}
 \frac{\partial\theta'}{\partial p}<0,
 \qquad \frac{\partial\theta_{-1}}{\partial p}>0.
\end{equation}
Both derivatives hold the current lifted angle fixed.  Thus increasing
momentum decreases the next angle and increases the preceding angle.

The generating function in these line coordinates
\citep{BialyMironov2017}
\begin{equation}\label{geo:generating}
 \cS(\theta,\theta')
 =2h\!\left(\frac{\theta+\theta'}2\right)
       \sin\!\left(\frac{\theta'-\theta}2\right)
\end{equation}
satisfies $p=-\cS_1(\theta,\theta')$ and
$p'=\cS_2(\theta,\theta')$.  Hence $T$ preserves $\dd p\wedge\dd\theta$.
The associated positive area measure on $\mathcal B$ is finite and has
full support.

The hypothesis supplies a uniform upper strip in these coordinates.
In fact, if the outgoing angle at the starting point is $\delta$, then
the starting normal angle is $\theta-\delta$ and
\begin{equation}\label{geo:starting-momentum}
 \begin{gathered}
 p=h(\theta-\delta)\cos\delta
                 +h'(\theta-\delta)\sin\delta,\\
 \frac{\partial p}{\partial\delta}
             =-r(\theta-\delta)\sin\delta<0.
 \end{gathered}
\end{equation}
Compactness and the inclusion of all sufficiently small positive
$\delta$ imply that, for some $\eps_g>0$,
\begin{equation}\label{geo:uniform-collar}
 \{(\theta,p):h(\theta)-\eps_g<p<h(\theta)\}\subset U.
\end{equation}

\subsection{A separation barrier for continuous leaves}

We record the elementary circle-dynamics facts that will also fix our
rotation convention.  If an increasing homeomorphism $G:\R\to\R$
commutes with translation by $2\pi$, its translation number satisfies
\begin{equation}\label{geo:translation-bound}
 \tau(G)=\lim_{j\to\infty}\frac{G^j(\alpha)-\alpha}{j},
 \qquad
 |G^j(\alpha)-\alpha-j\tau(G)|\le2\pi
 \quad (j\in\mathbb Z).
\end{equation}
For completeness, the displacements of each iterate oscillate by at most
$2\pi$, by monotonicity.  Their maxima are subadditive in the iteration
index, and their minima are superadditive.  Dividing by that index and
using repeated blocks of each fixed length gives a common limit bracketed
by the minimum and maximum divided by the index.  Their oscillation bound
then proves \eqref{geo:translation-bound}.  The negative-index case
follows by starting a positive iterate at $G^j(\alpha)$.
The same estimate, applied to a fixed large iterate, proves continuity of
$\tau$ under locally uniform convergence of lifts.  Also
$\tau(G^q)=q\tau(G)$.

One further consequence will be useful.  For integers $q\ge1$ and $l$,
\begin{equation}\label{geo:translation-sign}
 \begin{aligned}
 \tau(G)>2\pi l/q&\ \Longrightarrow\
              G^q(\alpha)>\alpha+2\pi l\quad\text{for every }\alpha,\\
 \tau(G)<2\pi l/q&\ \Longrightarrow\
              G^q(\alpha)<\alpha+2\pi l\quad\text{for every }\alpha.
 \end{aligned}
\end{equation}
A zero of the displacement difference would give a periodic orbit and
force translation number $2\pi l/q$.  In the absence of a zero, that
continuous periodic difference has a constant strict sign.  A uniformly
negative difference forces $q\tau(G)-2\pi l<0$ by iteration, and a
uniformly positive one forces the reverse inequality.  This proves
\eqref{geo:translation-sign}.

Use $b\in(0,1)$ for the transverse parameter in the given homeomorphism.
A constant-small-$\delta$ circle lies in $U$ and has winding one in the
phase cylinder. It is non-null in $U$ and, as an embedded essential circle
of the annulus, represents a primitive generator of $H_1(U;\mathbb Z)$.
Each product leaf represents that generator up to a common orientation
choice. The normal map of the positively curved boundary has degree one,
so the leaf's starting normal angle also has degree one. The line
angle is the starting normal angle plus the globally defined real angle
$\delta\in(0,\pi)$, so it has the same degree. Lift the foliation to
continuous parametrizations
\begin{equation}\label{geo:leaf-lifts}
 Z_b(\alpha)=(\theta_b(\alpha),p_b(\alpha)),\qquad
 Z_b(\alpha+2\pi)=Z_b(\alpha)+(2\pi,0).
\end{equation}
Each $Z_b$ parametrizes the full inverse image of its leaf injectively.
Individual invariance therefore defines a homeomorphism $G_b$ of $\R$ by
\[
 \widetilde T Z_b=Z_b\circ G_b.
\]
The deck relations for $Z_b$ and $\widetilde T$ give
$G_b(\alpha+2\pi)=G_b(\alpha)+2\pi$. A decreasing homeomorphism
cannot satisfy this relation, so $G_b$ is increasing.
The inverses of the lifted foliation charts are continuous,
so $G_b$ depends continuously on $b$.  Set $\tau(b)=\tau(G_b)$.
It is continuous by \eqref{geo:translation-bound}; because
$\theta_b(\alpha)-\alpha$ is bounded, it is also the limiting angular
advance per iterate.

Fix two leaves.  For distinct leaves use pairs $(\alpha,\beta)\in\R^2$;
for a single leaf use the connected strip
$\alpha<\beta\le\alpha+2\pi$.  Define
\begin{equation}\label{geo:separation}
 x_j=\theta_{b_2}(G_{b_2}^j(\beta))
                  -\theta_{b_1}(G_{b_1}^j(\alpha)).
\end{equation}
If $x_j=0$, the two momenta are distinct, by simplicity and disjointness
of leaves.  The period-separated pair in the single-leaf strip has
$x_j=2\pi$, so is not an exception.  The two inequalities in
\eqref{geo:twists} give
\begin{equation}\label{geo:opposite-signs}
 x_j=0\quad\Longrightarrow\quad x_{j-1}x_{j+1}<0.
\end{equation}
In particular the vector $(x_j,x_{j+1})$ never vanishes.

\begin{lemma}[Separation barrier]\label{geo:barrier}
For each fixed pair of leaves there is a continuous real argument
$A(\alpha,\beta)$ of $(x_0,x_1)$, periodic under simultaneous translation
of $\alpha,\beta$ by $2\pi$.  With
$A_j=A(G_{b_1}^j(\alpha),G_{b_2}^j(\beta))$, it satisfies
\begin{equation}\label{geo:argument-step}
 -\frac{3\pi}{2}<A_{j+1}-A_j<\frac\pi2,
\end{equation}
and, for every integer $m$,
\begin{equation}\label{geo:barrier-rule}
 A_j\le a_m:=\frac\pi2+\pi m
       \quad\Longrightarrow\quad A_{j+1}<a_m.
\end{equation}
If $x_i$ and $x_k$ have opposite strict signs for some $i<k$, then
$A_j$ cannot return arbitrarily close to $A_i$ as $j\to+\infty$.
\end{lemma}

\begin{proof}
The pair plane and pair strip are simply connected, so the nonvanishing
vector has a continuous real argument.  For a single leaf normalize it
to $\pi/4$ on $\beta=\alpha+2\pi$.  For distinct leaves, both entries
equal $\beta-\alpha+O(1)$, uniformly in the pair variables, so normalize
the argument to $(0,\pi/2)$ on the connected region where
$\beta-\alpha$ is sufficiently large.  Either normalization implies
simultaneous periodicity: the difference after a simultaneous translation
is a continuous integer multiple of $2\pi$ and vanishes in the
normalizing region.

Consecutive vectors have the form $(u,v),(v,w)$.  A positive
quarter-turn relation $(v,w)=c(-v,u)$, $c>0$, would force $v=0$
and $u,w$ to have the same sign, contrary to
\eqref{geo:opposite-signs}.  Thus $A_1-A_0$ avoids $\pi/2$ modulo
$2\pi$.  Connectedness, together with the edge normalization or the
large-separation normalization, puts it in $(-3\pi/2,\pi/2)$.
Applying this to the iterated pair proves \eqref{geo:argument-step}.

If $A_j\le a_m\le A_{j+1}$, that bound gives
\[
 a_m-\pi/2<A_j\le a_m\le A_{j+1}<a_m+\pi/2.
\]
On these two intervals $\sin A_j$ is nonzero with sign $(-1)^m$,
whereas $\cos A_{j+1}$ is either zero or has the opposite sign.  This
is impossible: both have the sign, and the same zero condition, of
$x_{j+1}$.  This proves \eqref{geo:barrier-rule}.

Finally, if $x_i\ne0$, then $A_i$ lies strictly between two consecutive
barriers.  All subsequent arguments stay below the upper one.  If
$x_k$ has the opposite sign, $A_k$ cannot remain in the interval
between those barriers, so lies strictly below its lower endpoint.
All later arguments remain below that endpoint, a positive distance
from $A_i$.
\end{proof}

\subsection{Graphs, recurrence, and rational rotation}

\begin{proposition}\label{prop:graph-collar}
Every leaf in $U$ is a graph $p=g_b(\theta)$.  After reversing $b$ if
necessary, $g_b$ is strictly decreasing in $b$, jointly continuous in
$(\theta,b)$, and converges uniformly to $h$ as $b\downarrow0$.
The induced angular lifts
\[
 F_b=\theta_b\circ G_b\circ\theta_b^{-1}
\]
are strictly ordered with $b$ and have positive translation numbers
tending to zero at that end.  Every sufficiently small positive rational
$l/q$ occurs on exactly one leaf, and on that entire leaf
\begin{equation}\label{geo:rational-periodicity}
 F_b^q(\theta)=\theta+2\pi l.
\end{equation}
Here rotation number means translation number divided by $2\pi$.
\end{proposition}

\begin{proof}
\emph{Periodic leaves are graphs.}
Suppose first that $G_b^q=\mathrm{id}+2\pi l$ on a whole leaf.
For its pair strip, simultaneous periodicity gives $A_q=A_0$.
If $x_0=0$, then $A_0$ is a barrier, and
\eqref{geo:barrier-rule} makes $A_q<A_0$, a contradiction.
Thus $x_0$ never vanishes on the connected strip.  It equals $2\pi$
on its upper edge, so is everywhere positive.  This says precisely
that $\theta_b$ is strictly increasing.

\emph{Constant-translation intervals also consist of graphs.}
Let $b$ be interior to an interval on which $\tau$ is constant.
If $\theta_b$ were not strictly increasing, its degree-one property
and continuity would give an ordered same-leaf pair with $x_0=0$.
By \eqref{geo:opposite-signs}, $x_{-1}$ and $x_1$ have opposite strict
signs.  Start instead with that preceding pair and perturb it slightly.
We obtain a nonempty open set $O$ of pairs on distinct leaves for which
$x_0x_2<0$.  We can take the two leaf parameters in disjoint relatively
compact intervals $I_1,I_2$ inside the constant-translation interval,
and take $|\beta-\alpha|$ bounded on $O$.

Consider the cover of
$H(\T\times I_1)\times H(\T\times I_2)$ obtained from
$(\R\times I_1)\times(\R\times I_2)$ by identifying only
\[
 (\alpha,b_1,\beta,b_2)
       \sim(\alpha+2\pi,b_1,\beta+2\pi,b_2).
\]
The pair dynamics $K$ induced by $(G_{b_1},G_{b_2})$ preserves the
measure lifted from product billiard area.  This measure is defined by
integrating counting measure on the covering fibers against product
area; $K$ bijects the fibers over the area-preserving base dynamics.
It has full support, and any region of bounded $|\beta-\alpha|$ has
finite measure, since it meets uniformly finitely many sheets over
each base point.  This construction requires no differentiability or
absolute continuity of $H$.

Equality of the two translation numbers and
\eqref{geo:translation-bound} imply, for every $j\in\mathbb Z$,
\begin{equation}\label{geo:relative-bound}
 \left|(G_{b_2}^j(\beta)-G_{b_1}^j(\alpha))
                         -(\beta-\alpha)\right|\le4\pi.
\end{equation}
Consequently the saturation $\bigcup_{j\in\mathbb Z}K^j(O)$ is an
invariant finite-measure set.  Poincar\'e recurrence gives recurrent
points in $O$, indeed almost every point there is recurrent.  One may
obtain recurrence to a point from the usual measurable-set version by
using a countable neighborhood base; the latter version follows because
the forward images of a set of nonreturners are disjoint in finite
measure.

For any such recurrent pair, the two leaf parameters remain fixed.
Thus continuity of the argument for this fixed pair of leaves and
simultaneous periodicity give $A_{j_k}\to A_0$ along a sequence
$j_k\to\infty$.  This contradicts Lemma~\ref{geo:barrier}, because
$x_0x_2<0$.  Hence $\theta_b$ is strictly increasing throughout every
constant-translation interval.

\emph{Graph leaves are dense, and then every leaf is a graph.}
Take any open parameter interval $J$.  If it contains a subinterval
where $\tau$ is constant, the preceding argument supplies graph leaves.
Otherwise choose a rational multiple $r_*=2\pi l/q$ strictly between
two values of the continuous function $\tau$ on $J$.  Travel along
the parameter segment from the smaller value to the larger one, and
let $c$ be the first position approached by values $\tau>r_*$.
Continuity gives $\tau(c)=r_*$.  Values strictly below $r_*$ also
approach $c$: otherwise $\tau=r_*$ on an interval immediately before
$c$, contrary to our case assumption.  Taking limits in
\eqref{geo:translation-sign} from these two strict sides gives
$G_c^q=\mathrm{id}+2\pi l$.  The periodic-leaf argument therefore
supplies a graph in $J$.  Notice that no monotonicity of $\tau$ has
been used.

Graph parameters are thus dense.  Continuity of the lifted foliation
makes every $\theta_b$ nondecreasing.  If a same-leaf ordered pair had
$x_0=0$, its preceding and following pairs would still be ordered,
because $G_b$ is increasing.  Nondecrease would give
$x_{-1},x_1\ge0$, contradicting \eqref{geo:opposite-signs}.
Every $\theta_b$ is therefore strictly increasing, proving the graph
property.

\emph{Ordering and the grazing end.}
Inverting the increasing angular parametrizations shows that $g_b$ is
jointly continuous.  Distinct graphs are disjoint, so their order cannot
change with $\theta$.  For each fixed angle, the continuous injective
function $b\mapsto g_b(\theta)$ is strictly monotone, in the same
direction at every angle.  Reverse $b$ to make this direction
decreasing.  The upper-strip inclusion \eqref{geo:uniform-collar}
then forces $g_b(\theta)\to h(\theta)$ as $b\downarrow0$ at every
angle.  Monotonicity and compactness upgrade this to uniform convergence.
Formula \eqref{geo:reflection} now shows that $F_b(\theta)-\theta$ is
positive and tends uniformly to zero at that end.  Its positive minimum
on each fixed graph gives $\tau(b)>0$, and $\tau(b)\to0$ follows.

The forward twist and decreasing graph order give $F_{b_1}<F_{b_2}$
when $b_1<b_2$.  Since the maps are increasing, all their positive
iterates have the same strict order.  Passing to translation numbers
shows that $\tau(b)$ is nondecreasing.

\emph{Rational plateaus are impossible.}
If $\tau(b)=2\pi l/q$ on an open parameter interval $I$, almost every
point in its finite-area invariant annulus is recurrent for $T^q$.
On each leaf, $F_b^q-2\pi l$ has zero translation number and hence has
fixed points: without a fixed point its periodic displacement would
have a uniform strict sign.  Every nonfixed point moves monotonically
in a component of the complement of the fixed set towards an endpoint,
so is not recurrent on the circle.  Thus $T^q$ is the identity almost
everywhere in this annulus, and therefore everywhere, by continuity and
full support of area.  Its lift on each graph is necessarily
$F_b^q=\mathrm{id}+2\pi l$, with the integer shift determined by the
translation number.  This contradicts strict ordering of the iterates
on two distinct graphs.

The continuity, monotonicity, positivity, and zero limiting value of
$\tau$ now show that every sufficiently small positive rational
$l/q$ occurs.  Its parameter is unique, since two parameters with that
value would give a plateau.  Values on either side are strictly below
and above $2\pi l/q$; taking limits in
\eqref{geo:translation-sign} proves
\eqref{geo:rational-periodicity} on the entire graph.
\end{proof}

\subsection{Closing stationary chains}

Proposition~\ref{prop:graph-collar} supplies the phase conclusion used in
\cite[Lemma~2.2]{NB}.  We isolate an immediate variational consequence
because it will be used twice: a fixed-endpoint stationary chain closes
in momentum as well as in angle. This rational-graph/action connection is
related to the periodic-graph regularity studied by
\citet[Appendix, Proposition~A.1]{FierobeSorrentino2026}; our argument
starts with the full continuous foliation and proves the graph and closure
properties before any analytic regularity is known.

\begin{lemma}\label{lem:stationary-closure}
There is $\kappa>0$ with the following property.  Let $q\ge1$ and $l$
be integers, and let $\theta_0,\ldots,\theta_q$ be a real chain such that
\[
 0<\theta_{j+1}-\theta_j<\kappa,\qquad
 \theta_q=\theta_0+2\pi l,
\]
and suppose it is stationary at all interior nodes:
\begin{equation}\label{geo:interior-stationarity}
 \cS_2(\theta_{j-1},\theta_j)
             +\cS_1(\theta_j,\theta_{j+1})=0,
 \qquad 1\le j<q.
\end{equation}
Define the momenta, including the final one, by
\begin{equation}\label{geo:chain-momenta}
 p_j=-\cS_1(\theta_j,\theta_{j+1})\quad(0\le j<q),
 \qquad p_q=\cS_2(\theta_{q-1},\theta_q).
\end{equation}
Then $\widetilde T(\theta_j,p_j)=(\theta_{j+1},p_{j+1})$ for every
$0\le j<q$, and $p_q=p_0$.  The chain therefore extends to a periodic
orbit of rotation $l/q$, on the unique graph with that rotation.
The action $\sum_{j=0}^{q-1}\cS(\theta_j,\theta_{j+1})$ is constant
along any $C^1$ family of such chains with fixed $q,l$.
\end{lemma}

\begin{proof}
Choose $\kappa<2\pi$ small enough that
\eqref{geo:reflection} and \eqref{geo:uniform-collar} place every line
with angular step in $(0,\kappa)$ in $U$.  Decrease it further so that
Proposition~\ref{prop:graph-collar} applies to every positive rational
smaller than $\kappa/(2\pi)$.  The generating relation and
\eqref{geo:interior-stationarity} give every stated transition; for
the last transition use precisely the definition of $p_q$ in
\eqref{geo:chain-momenta}.  No endpoint stationarity has been assumed.

The initial state belongs to an individually invariant graph in $U$,
so every state of the chain belongs to that same graph.  The equality
$\theta_q=\theta_0+2\pi l$ therefore forces $p_q=p_0$.
Repeating the orbit now shows that its translation number is
$2\pi l/q$.  Since $0<l/q<\kappa/(2\pi)$, uniqueness of the rational
graph applies.  In particular, all chains with the same $q,l$ lie on
the same graph, regardless of their starting angles.

Finally, differentiating the action along a $C^1$ family cancels all
interior terms by \eqref{geo:interior-stationarity}.  The remaining
terms are $-p_0\dot\theta_0+p_q\dot\theta_q=0$, because the endpoint
advance is fixed and the momenta agree.
\end{proof}
